\chapter{Chapter05 — Peres}
%%% generated by the blueprint pipeline from FABL.Chapter05.Peres
\section{Overview}

Declarations of the Lean module \texttt{FABL.Chapter05.Peres}.

\section{Declarations}

\begin{definition}[Linear threshold class]
\lean{FABL.linearThresholdClass}
\statementsource{Analysis-of-Boolean-Functions-by-Ryan-ODonnell}{
  pdf-ffbe9c7b3306-p044-b005,
  pdf-ffbe9c7b3306-p113-b003
}
\label{FABL.linearThresholdClass}
\leanok
\uses{FABL.BooleanClass, FABL.BooleanFunction, FABL.IsLinearThreshold}
The linear threshold class assigns to each natural number $n$ the set of Boolean
functions on $n$ coordinates that are linear threshold functions.
\end{definition}

\begin{theorem}[Sign-flip invariance of linear threshold functions]
\lean{FABL.linearThresholdClass_closedUnderNegatingInputVariables}
\label{FABL.linearThresholdClass_closedUnderNegatingInputVariables}
\leanok
\uses{FABL.IsClosedUnderNegatingInputVariables, FABL.IsLinearThreshold, FABL.linearThresholdClass, FABL.negateInputVariables, FABL.signValue, FABL.thresholdSign}
\difficulty{4}
For every dimension $n$, every sign vector $z\in\{\pm1\}^n$, and every linear threshold
function $f:\{\pm1\}^n\to\{\pm1\}$, the function $x\mapsto f(z_1x_1,\ldots,z_nx_n)$ is
also a linear threshold function.
\end{theorem}

\begin{theorem}[Closure under identifying input variables]
\lean{FABL.linearThresholdClass_closedUnderIdentifyingInputVariables}
\proofsource{Analysis-of-Boolean-Functions-by-Ryan-ODonnell}{pdf-ffbe9c7b3306-p132-b002}
\label{FABL.linearThresholdClass_closedUnderIdentifyingInputVariables}
\leanok
\uses{FABL.IsClosedUnderIdentifyingInputVariables, FABL.IsLinearThreshold, FABL.identifyInputVariables, FABL.linearThresholdClass, FABL.signValue, FABL.thresholdSign}
\difficulty{4}
For all natural numbers $n,m$, every map $\pi:\{1,\ldots,n\}\to\{1,\ldots,m\}$, and
every linear threshold function $f:\{-1,1\}^n\to\{-1,1\}$, the function $w\mapsto
f(w_{\pi(1)},\ldots,w_{\pi(n)})$ is also a linear threshold function. Here a linear
threshold function has the form $f(x)=1$ when $a_0+\sum_{i=1}^n a_ix_i\ge 0$ and
$f(x)=-1$ otherwise, for some real coefficients $a_0,a_1,\ldots,a_n$.
\end{theorem}

\begin{theorem}[Influence bound for linear threshold functions]
\lean{FABL.totalInfluence_toReal_le_sqrt_of_mem_linearThresholdClass}
\proofsource{Analysis-of-Boolean-Functions-by-Ryan-ODonnell}{
  pdf-ffbe9c7b3306-p063-b002,
  pdf-ffbe9c7b3306-p132-b002
}
\label{FABL.totalInfluence_toReal_le_sqrt_of_mem_linearThresholdClass}
\leanok
\uses{FABL.BooleanFunction, FABL.BooleanFunction.toReal, FABL.isUnate_of_isLinearThreshold, FABL.linearThresholdClass, FABL.totalInfluence, FABL.totalInfluence_toReal_le_sqrt_card_of_unate}
\difficulty{1}
Let $r$ be a positive integer and let $f:\{-1,1\}^r\to\{-1,1\}$ be a linear
threshold function: there are real weights $a_1,\ldots,a_r$ and a real threshold offset
$a_0$ such that $f(x)$ is the sign of $a_0+\sum_{i=1}^r a_i x_i$, with a nonnegative
input sent to $1$ and a negative input sent to $-1$. Then the total influence of $f$ is
at most $\sqrt{r}$.
\end{theorem}

\begin{theorem}[Square-root noise bound for threshold functions]
\lean{FABL.noiseSensitivity_le_sqrt_inverseNoiseFloor_of_isLinearThreshold}
\proofsource{Analysis-of-Boolean-Functions-by-Ryan-ODonnell}{
  pdf-ffbe9c7b3306-p131-b002,
  pdf-ffbe9c7b3306-p131-b003,
  pdf-ffbe9c7b3306-p132-b002
}
\label{FABL.noiseSensitivity_le_sqrt_inverseNoiseFloor_of_isLinearThreshold}
\leanok
\uses{FABL.BooleanFunction, FABL.IsLinearThreshold, FABL.PositiveHalfNoiseParameter, FABL.inverseNoiseFloor, FABL.linearThresholdClass, FABL.linearThresholdClass_closedUnderIdentifyingInputVariables, FABL.linearThresholdClass_closedUnderNegatingInputVariables, FABL.noiseSensitivity, FABL.noiseSensitivity_le_inverseNoiseFloor_totalInfluenceBound, FABL.totalInfluence_toReal_le_sqrt_of_mem_linearThresholdClass}
\difficulty{3}
Let $n$ be a nonnegative integer and let $f:\{-1,1\}^n\to\{-1,1\}$ be a linear threshold
function: there are real weights $a_1,\ldots,a_n$ and a real offset $a_0$ such that
$f(x)$ is $1$ when $a_0+\sum_{i=1}^n a_i x_i\geq0$ and is $-1$ otherwise. For
$0<\delta\leq\tfrac12$, the probability that $f$ changes when each coordinate of a
uniformly chosen input is independently flipped with probability $\delta$ is at most
$\sqrt{1/m}$, where $m=\lfloor1/\delta\rfloor$.
\end{theorem}

\begin{theorem}[Inverse Noise Floor Bound]
\lean{FABL.sqrt_inverseNoiseFloor_le_sqrt_three_halves_mul_sqrt}
\label{FABL.sqrt_inverseNoiseFloor_le_sqrt_three_halves_mul_sqrt}
\leanok
\uses{FABL.PositiveHalfNoiseParameter, FABL.inverseNoiseFloor}
\difficulty{5}
Let $\delta$ be a real number with $0<\delta\le \tfrac12$, and let $N=\lfloor
1/\delta\rfloor$, the greatest integer at most $1/\delta$. Then $\sqrt{1/N}\le
\sqrt{3/2}\,\sqrt{\delta}$.
\end{theorem}

\begin{theorem}[Bounds for a reciprocal floor]
\lean{FABL.reciprocal_natFloor_one_div_bounds}
\label{FABL.reciprocal_natFloor_one_div_bounds}
\leanok
\difficulty{5}
For every real number $\delta$ with $0<\delta\le 1/2$, let $m=\lfloor 1/\delta\rfloor$,
where the floor is the greatest integer not exceeding $1/\delta$. Then $\delta\le 1/m\le
2\delta$, and $1/m-\delta\le 2\delta^2$.
\end{theorem}

\begin{theorem}[Floor Approximation Error Bound]
\lean{FABL.sqrt_one_div_natFloor_one_div_sub_sqrt_isBigO}
\label{FABL.sqrt_one_div_natFloor_one_div_sub_sqrt_isBigO}
\leanok
\uses{FABL.reciprocal_natFloor_one_div_bounds, FABL.sqrt_cube_eq_rpow_three_halves}
\difficulty{6}
As $\delta$ tends to $0$ through nonnegative real values, $\sqrt{1/\lfloor
1/\delta\rfloor}-\sqrt{\delta}$ is $O(\delta^{3/2})$, where $\lfloor 1/\delta\rfloor$ is
the greatest integer at most $1/\delta$.
\end{theorem}

\begin{theorem}[Noise sensitivity bound]
\lean{FABL.peresNoiseSensitivityBound}
\proofsource{Analysis-of-Boolean-Functions-by-Ryan-ODonnell}{
  pdf-ffbe9c7b3306-p131-b003,
  pdf-ffbe9c7b3306-p132-b002,
  pdf-ffbe9c7b3306-p130-b004
}
\label{FABL.peresNoiseSensitivityBound}
\leanok
\uses{FABL.BooleanFunction, FABL.IsLinearThreshold, FABL.PositiveHalfNoiseParameter, FABL.noiseSensitivity, FABL.noiseSensitivity_le_sqrt_inverseNoiseFloor_of_isLinearThreshold, FABL.sqrt_inverseNoiseFloor_le_sqrt_three_halves_mul_sqrt}
\difficulty{1}
For every natural number $n$, every linear threshold function $f:\{\pm1\}^n\to\{\pm1\}$,
and every noise parameter $\delta$ with $0<\delta\le 1/2$, the noise sensitivity of $f$
at parameter $\delta$ is at most $\sqrt{3/2}\,\sqrt{\delta}$.
\end{theorem}

\begin{theorem}[Total influence majority le main add four div sqrt]
\lean{FABL.totalInfluence_majority_le_main_add_four_div_sqrt}
\statementsource{Analysis-of-Boolean-Functions-by-Ryan-ODonnell}{
  pdf-ffbe9c7b3306-p051-b001,
  pdf-ffbe9c7b3306-p063-b001,
  pdf-ffbe9c7b3306-p119-b001,
  pdf-ffbe9c7b3306-p119-b005,
  pdf-ffbe9c7b3306-p120-b001
}
\label{FABL.totalInfluence_majority_le_main_add_four_div_sqrt}
\leanok
\uses{FABL.BooleanFunction.toReal, FABL.abs_evenMajorityTotalInfluenceError_le, FABL.evenMajorityTotalInfluenceError, FABL.majority, FABL.majority_isMajorityFunction, FABL.oddMajorityTotalInfluenceError, FABL.oddMajorityTotalInfluenceError_mem_Icc, FABL.totalInfluence, FABL.totalInfluence_evenMajority_eq_main_add_error, FABL.totalInfluence_majority_odd_eq_main_add_error}
\difficulty{5}
% dataset-pipeline: stub -- prose pending build_informal_statements.py
Statement pending regeneration.
\end{theorem}

\begin{theorem}[Square-root quotient identity]
\lean{FABL.main_add_four_div_sqrt_quotient_eq}
\label{FABL.main_add_four_div_sqrt_quotient_eq}
\leanok
\uses{FABL.sqrt_cube_eq_rpow_three_halves}
\difficulty{4}
For every real number $m>0$,
\[\frac{\sqrt{2/\pi}\,\sqrt m+4/\sqrt m}{m}=\sqrt{2/\pi}\,\sqrt{1/m}+4(1/m)^{3/2}.\]
\end{theorem}

\begin{definition}[Peres majority upper bound]
\lean{FABL.peresMajorityUpperBound}
\statementsource{Analysis-of-Boolean-Functions-by-Ryan-ODonnell}{
  pdf-ffbe9c7b3306-p119-b001,
  pdf-ffbe9c7b3306-p119-b005,
  pdf-ffbe9c7b3306-p120-b001,
  pdf-ffbe9c7b3306-p130-b004,
  pdf-ffbe9c7b3306-p131-b002,
  pdf-ffbe9c7b3306-p131-b003,
  pdf-ffbe9c7b3306-p132-b001,
  pdf-ffbe9c7b3306-p132-b002
}
\label{FABL.peresMajorityUpperBound}
\leanok
For a real number $\delta$, define the Peres majority upper bound to be
$\sqrt{2/\pi}\,\sqrt{1/\lfloor 1/\delta\rfloor}+4(1/\lfloor 1/\delta\rfloor)^{3/2}$,
where the floor denotes the greatest natural number at most its argument, viewed as a
real number in the formula.
\end{definition}

\begin{theorem}[Noise sensitivity le majority influence ratio of is linear threshold]
\lean{FABL.noiseSensitivity_le_majorityInfluenceRatio_of_isLinearThreshold}
\statementsource{Analysis-of-Boolean-Functions-by-Ryan-ODonnell}{
  pdf-ffbe9c7b3306-p131-b002,
  pdf-ffbe9c7b3306-p131-b003,
  pdf-ffbe9c7b3306-p132-b001
}
\label{FABL.noiseSensitivity_le_majorityInfluenceRatio_of_isLinearThreshold}
\leanok
\uses{FABL.BooleanFunction, FABL.BooleanFunction.toReal, FABL.IsLinearThreshold, FABL.PositiveHalfNoiseParameter, FABL.inverseNoiseFloor, FABL.isUnate_of_isLinearThreshold, FABL.linearThresholdClass, FABL.linearThresholdClass_closedUnderIdentifyingInputVariables, FABL.linearThresholdClass_closedUnderNegatingInputVariables, FABL.majority, FABL.noiseSensitivity, FABL.noiseSensitivity_le_inverseNoiseFloor_totalInfluenceBound, FABL.totalInfluence, FABL.totalInfluence_toReal_le_majority_of_unate}
\difficulty{3}
% dataset-pipeline: stub -- prose pending build_informal_statements.py
Statement pending regeneration.
\end{theorem}

\begin{theorem}[Noise sensitivity bound for threshold functions]
\lean{FABL.noiseSensitivity_le_peresMajorityUpperBound}
\proofsource{Analysis-of-Boolean-Functions-by-Ryan-ODonnell}{
  pdf-ffbe9c7b3306-p131-b003,
  pdf-ffbe9c7b3306-p132-b002,
  pdf-ffbe9c7b3306-p130-b004
}
\label{FABL.noiseSensitivity_le_peresMajorityUpperBound}
\leanok
\uses{FABL.BooleanFunction, FABL.BooleanFunction.toReal, FABL.IsLinearThreshold, FABL.PositiveHalfNoiseParameter, FABL.inverseNoiseFloor, FABL.main_add_four_div_sqrt_quotient_eq, FABL.majority, FABL.noiseSensitivity, FABL.noiseSensitivity_le_majorityInfluenceRatio_of_isLinearThreshold, FABL.peresMajorityUpperBound, FABL.totalInfluence, FABL.totalInfluence_majority_le_main_add_four_div_sqrt}
\difficulty{4}
Let $n$ be a nonnegative integer and let $f:\{-1,1\}^n\to\{-1,1\}$ be a linear threshold
function: there are real weights $a_1,\ldots,a_n$ and a real offset $a_0$ such that
$f(x)$ is $1$ when $a_0+\sum_{i=1}^n a_i x_i\geq0$ and is $-1$ otherwise. For
$0<\delta\leq\tfrac12$, the probability that $f$ changes when each coordinate of a
uniformly chosen input is independently flipped with probability $\delta$ is at most
$\sqrt{2/\pi}\sqrt{1/\lfloor1/\delta\rfloor}+4(1/\lfloor1/\delta\rfloor)^{3/2}$.
\end{theorem}

\begin{theorem}[Reciprocal floor asymptotic bound]
\lean{FABL.one_div_natFloor_one_div_isBigO}
\label{FABL.one_div_natFloor_one_div_isBigO}
\leanok
\uses{FABL.reciprocal_natFloor_one_div_bounds}
\difficulty{4}
As $\delta$ tends to $0$ through positive real values, the reciprocal of $\lfloor
1/\delta \rfloor$ is $O(\delta)$; that is, the function $\delta \mapsto 1/\lfloor
1/\delta \rfloor$ is bounded in magnitude by a constant multiple of $\delta$ on some
right neighborhood of zero.
\end{theorem}

\begin{theorem}[Asymptotic remainder bound]
\lean{FABL.peresMajorityUpperBound_sub_main_isBigO}
\label{FABL.peresMajorityUpperBound_sub_main_isBigO}
\leanok
\uses{FABL.one_div_natFloor_one_div_isBigO, FABL.peresMajorityUpperBound, FABL.sqrt_one_div_natFloor_one_div_sub_sqrt_isBigO}
\difficulty{4}
As $\delta$ tends to $0$ through nonnegative values, the difference between the Peres
majority upper bound at $\delta$ and $\sqrt{2/\pi}\,\sqrt{\delta}$ is $O(\delta^{3/2})$.
\end{theorem}
